\section{Bases teóricas}
\label{sec:bases_teoricas}

El sistema se apoya en ideas ya probadas de informática, ingeniería de software y gestión escolar:

\subsection{Arquitectura Cliente-Servidor y Patrón MVC}

De acuerdo con \textcite{tanenbaum2021}, la arquitectura cliente-servidor constituye el modelo estructural predominante en internet, donde el navegador (cliente) interactúa vía HTTP con el servidor Apache/PHP y la base de datos relacional. Sobre esta infraestructura, el proyecto adopta el patrón arquitectónico Modelo--Vista--Controlador (MVC) \parencite{pressman2020}, que desacopla responsabilidades: la \textbf{Vista} (HTML5/CSS3/JavaScript y plantillas PHP) gestiona la presentación y validación reactiva en cliente; el \textbf{Controlador} (PHP 8.2/Apache) centraliza el enrutamiento, sesiones y reglas de negocio (recovery C.I.+fecha, matrícula y paralelos); y el \textbf{Modelo} (clases de dominio PDO--PostgreSQL 16) encapsula la lógica de persistencia ACID. Esta separación garantiza modularidad, mantenibilidad y mínimo consumo de ancho de banda en redes móviles de áreas intermedias.

\subsection{Ingeniería de Software Web y Usabilidad Centrada en el Usuario}

\textcite{pressman2020} describen la ingeniería de software web (WebE) como la forma ordenada de crear aplicaciones seguras y que se puedan mantener en la red. En un colegio, donde cada estudiante maneja la tecnología de manera distinta, la facilidad de uso es clave. Por eso se buscó que los formularios sean simples, avisen al instante si falta un dato y respondan rápido, para que inscribirse no genere dudas.

\subsection{Seguridad Informática, Autenticación y Recuperación de Credenciales}

La protección de los datos y el control de acceso siguen los principios de \textcite{stallings2021}: confidencialidad, integridad y disponibilidad. Las contraseñas se guardan cifradas con funciones hash en el servidor. La idea de autoservicio de recuperación (Self-Service Password Reset) plantea que el propio usuario pueda restablecer su cuenta comprobando datos que solo él conoce, como el número de carnet y la fecha de nacimiento, sin tener que ir a ventanilla ni quedar bloqueado.

\subsection{Sistemas de Información y Gestión Documental Escolar}

En la práctica escolar, inscribirse no termina al reservar el cupo. Hace falta comprobar los papeles que pide la norma. Por eso el sistema incluye un módulo donde se anotan y revisan las libretas, las cédulas y los certificados que el estudiante presenta en físico, de modo que los expedientes dejan de estar dispersos y quedan registrados de forma ordenada y rastreable.

\subsection{Modelado Orientado a Objetos y Persistencia Transaccional ACID en PostgreSQL}

Según \textcite{sommerville2020}, usar UML permite dibujar con claridad las clases y el flujo del sistema antes de programar. La guarda segura de los cupos se logra con PostgreSQL \parencite{postgresql2025}, que cumple las propiedades ACID (Atomicidad, Consistencia, Aislamiento y Durabilidad) y evita que dos postulantes ocupen el mismo lugar o que una asignación quede a medias.
